% Half-Life Calculations Worksheet
\documentclass[12pt]{article}
\usepackage[margin=1in]{geometry}
\usepackage{multicol}
\usepackage{siunitx}
\usepackage{amsmath}
\setlength{\parindent}{0pt}

\begin{document}

\begin{center}
{\LARGE Worksheet: Half-Life Calculations}\\
{\large Fractions and Mass Remaining}
\end{center}

\vspace{0.5em}
Name: \rule{0.6\textwidth}{0.4pt} \hfill Date: \rule{0.2\textwidth}{0.4pt}

% ---------------------- Section 1 ----------------------
\section*{Section A — Fraction Remaining}
State the fraction of the original sample remaining after the given number of half-lives.

\begin{enumerate}
\item Strontium-90 undergoes beta decay. After 4 half-lives, what fraction remains?
\item Iodine-131 undergoes beta decay. After 3 half-lives, what fraction remains?
\item Phosphorus-32 undergoes beta decay. After 6 half-lives, what fraction remains?
\item Carbon-14 undergoes beta decay. After 8 half-lives, what fraction remains?
\item Radon-222 undergoes alpha decay. After 2 half-lives, what fraction remains?
\item Polonium-210 undergoes alpha decay. After 5 half-lives, what fraction remains?
\item Thorium-232 undergoes alpha decay. After 7 half-lives, what fraction remains?
\item Uranium-238 undergoes alpha decay. After 1 half-life, what fraction remains?
\item Radium-226 undergoes alpha decay. After 6 half-lives, what fraction remains?
\item Cobalt-60 undergoes beta decay. After 2 half-lives, what fraction remains?
\end{enumerate}

% ---------------------- Section 2 ----------------------
\section*{Section B — Mass Remaining from Half-Lives}
Find the mass remaining given the initial mass and a number of half-lives.

\begin{enumerate}
\setcounter{enumi}{10}
\item Strontium-90: An initial sample of \(\SI{240}{g}\) decays for 4 half-lives. What mass remains?
\item Iodine-131: An initial sample of \(\SI{300}{g}\) decays for 2 half-lives. What mass remains?
\item Technetium-99m: An initial sample of \(\SI{256}{g}\) decays for 5 half-lives. What mass remains?
\item Phosphorus-32: An initial sample of \(\SI{96}{g}\) decays for 3 half-lives. What mass remains?
\item Radon-222: An initial sample of \(\SI{160}{g}\) decays for 5 half-lives. What mass remains?
\item Fluorine-18: An initial sample of \(\SI{200}{g}\) decays for 3 half-lives. What mass remains?
\item Sodium-24: An initial sample of \(\SI{384}{g}\) decays for 7 half-lives. What mass remains?
\item Chromium-51: An initial sample of \(\SI{72}{g}\) decays for 2 half-lives. What mass remains?
\item Molybdenum-99: An initial sample of \(\SI{1024}{g}\) decays for 8 half-lives. What mass remains?
\item Iodine-123: An initial sample of \(\SI{350}{g}\) decays for 1 half-life. What mass remains?
\end{enumerate}

% ---------------------- Section 3 ----------------------
\section*{Section C — Mass from Time}
Use the given half-life and elapsed time to determine the number of half-lives, then find the mass remaining.

\begin{enumerate}
\setcounter{enumi}{20}
\item Iodine-131 has a half-life of \(\SI{8}{days}\). A \(\SI{240}{g}\) sample decays for \(\SI{32}{days}\). Find the mass remaining.
\item Technetium-99m has a half-life of \(\SI{6}{hours}\). A \(\SI{400}{g}\) sample decays for \(\SI{18}{hours}\). Find the mass remaining.
\item Fluorine-18 has a half-life of \(\SI{110}{minutes}\). A \(\SI{160}{g}\) sample decays for \(\SI{440}{minutes}\). Find the mass remaining.
\item Sodium-24 has a half-life of \(\SI{15}{hours}\). A \(\SI{256}{g}\) sample decays for \(\SI{60}{hours}\). Find the mass remaining.
\item Phosphorus-32 has a half-life of \(\SI{14}{days}\). A \(\SI{300}{g}\) sample decays for \(\SI{28}{days}\). Find the mass remaining.
\item Chromium-51 has a half-life of \(\SI{28}{days}\). A \(\SI{128}{g}\) sample decays for \(\SI{84}{days}\). Find the mass remaining.
\item Radon-222 has a half-life of \(\SI{3.8}{days}\). A \(\SI{96}{g}\) sample decays for \(\SI{7.6}{days}\). Find the mass remaining.
\item Molybdenum-99 has a half-life of \(\SI{66}{hours}\). A \(\SI{512}{g}\) sample decays for \(\SI{198}{hours}\). Find the mass remaining.
\item Nitrogen-13 has a half-life of \(\SI{10}{minutes}\). A \(\SI{320}{g}\) sample decays for \(\SI{50}{minutes}\). Find the mass remaining.
\item Iodine-123 has a half-life of \(\SI{13}{hours}\). A \(\SI{400}{g}\) sample decays for \(\SI{52}{hours}\). Find the mass remaining.
\end{enumerate}

\newpage
\begin{center}
{\LARGE Answer Key}
\end{center}

\textbf{Section A}
\begin{enumerate}
\item \(\left(\tfrac{1}{2}\right)^4=\tfrac{1}{16}\)
\item \(\left(\tfrac{1}{2}\right)^3=\tfrac{1}{8}\)
\item \(\left(\tfrac{1}{2}\right)^6=\tfrac{1}{64}\)
\item \(\left(\tfrac{1}{2}\right)^8=\tfrac{1}{256}\)
\item \(\left(\tfrac{1}{2}\right)^2=\tfrac{1}{4}\)
\item \(\left(\tfrac{1}{2}\right)^5=\tfrac{1}{32}\)
\item \(\left(\tfrac{1}{2}\right)^7=\tfrac{1}{128}\)
\item \(\left(\tfrac{1}{2}\right)^1=\tfrac{1}{2}\)
\item \(\left(\tfrac{1}{2}\right)^6=\tfrac{1}{64}\)
\item \(\left(\tfrac{1}{2}\right)^2=\tfrac{1}{4}\)
\end{enumerate}

\textbf{Section B}
\begin{enumerate}
\setcounter{enumi}{10}
\item \(240\times\tfrac{1}{16}=\SI{15}{g}\)
\item \(300\times\tfrac{1}{4}=\SI{75}{g}\)
\item \(256\times\tfrac{1}{32}=\SI{8}{g}\)
\item \(96\times\tfrac{1}{8}=\SI{12}{g}\)
\item \(160\times\tfrac{1}{32}=\SI{5}{g}\)
\item \(200\times\tfrac{1}{8}=\SI{25}{g}\)
\item \(384\times\tfrac{1}{128}=\SI{3}{g}\)
\item \(72\times\tfrac{1}{4}=\SI{18}{g}\)
\item \(1024\times\tfrac{1}{256}=\SI{4}{g}\)
\item \(350\times\tfrac{1}{2}=\SI{175}{g}\)
\end{enumerate}

\textbf{Section C}
\begin{enumerate}
\setcounter{enumi}{20}
\item \(n=\tfrac{32}{8}=4\Rightarrow 240\times\tfrac{1}{16}=\SI{15}{g}\)
\item \(n=\tfrac{18}{6}=3\Rightarrow 400\times\tfrac{1}{8}=\SI{50}{g}\)
\item \(n=\tfrac{440}{110}=4\Rightarrow 160\times\tfrac{1}{16}=\SI{10}{g}\)
\item \(n=\tfrac{60}{15}=4\Rightarrow 256\times\tfrac{1}{16}=\SI{16}{g}\)
\item \(n=\tfrac{28}{14}=2\Rightarrow 300\times\tfrac{1}{4}=\SI{75}{g}\)
\item \(n=\tfrac{84}{28}=3\Rightarrow 128\times\tfrac{1}{8}=\SI{16}{g}\)
\item \(n=\tfrac{7.6}{3.8}=2\Rightarrow 96\times\tfrac{1}{4}=\SI{24}{g}\)
\item \(n=\tfrac{198}{66}=3\Rightarrow 512\times\tfrac{1}{8}=\SI{64}{g}\)
\item \(n=\tfrac{50}{10}=5\Rightarrow 320\times\tfrac{1}{32}=\SI{10}{g}\)
\item \(n=\tfrac{52}{13}=4\Rightarrow 400\times\tfrac{1}{16}=\SI{25}{g}\)
\end{enumerate}

\end{document}

